\section{Results}
\label{sec:results}

Figures are for the validation year (2014) unless stated otherwise;
held-out test-year (2015) confirmation, evaluated once after every
modelling and evaluation choice was frozen, is reported alongside each
finding. Table~\ref{tab:T2} gives the array-11 configurations
underlying both effects in this section; arrays 12 and 17 are reported
in the text below.

% T2_protocol.tex - protocol-inflation configurations (array11), booktabs fragment.
% Source: paper/overleaf/tables/T4_protocol.tex, content unchanged,
% renamed T2 (second table in this paper) and relabeled tab:T2.
% Needs \usepackage{booktabs}.
% SCOPE: array11 only - see paper/overleaf/sections/04_methodology.tex
% for why (this array's evaluable sample is representative and
% avoids array17's outage-adjusted count).
\begin{table*}[t]
  \centering
  \small
  \caption{Protocol configurations C1-C6 (defined in the text above),
  array 11, lagged regime, validation year 2014: sample count,
  normalised RMSE, skill score, and reference forecast at each
  horizon. Arrays 12 and 17 show the same pattern
  (Section~\ref{sec:results}).}
  \label{tab:T2}
  \begin{tabular}{cl rrrl}
    \toprule
    $h$ & Config & $n$ & nRMSE (\%) & skill & Reference \\
    \midrule
    1 & C1 & 3773 & 6.36 & +0.200 & convex \\
     & C2 & 3773 & 6.36 & +0.252 & persistence \\
     & C3 & 4255 & 6.00 & +0.207 & convex \\
     & C4 & 8694 & 4.20 & +0.261 & persistence \\
     & C5 & 3773 & 6.36 & -- & -- \\
     & C6 & 8694 & 4.20 & -- & -- \\
    \addlinespace
    3 & C1 & 3762 & 8.97 & +0.276 & convex \\
     & C2 & 3762 & 8.97 & +0.526 & persistence \\
     & C3 & 4247 & 8.44 & +0.276 & convex \\
     & C4 & 8694 & 5.90 & +0.528 & persistence \\
     & C5 & 3762 & 8.97 & -- & -- \\
     & C6 & 8694 & 5.90 & -- & -- \\
    \addlinespace
    6 & C1 & 3756 & 10.71 & +0.194 & convex \\
     & C2 & 3756 & 10.71 & +0.652 & persistence \\
     & C3 & 4245 & 10.07 & +0.194 & convex \\
     & C4 & 8694 & 7.04 & +0.652 & persistence \\
     & C5 & 3756 & 10.71 & -- & -- \\
     & C6 & 8694 & 7.04 & -- & -- \\
    \bottomrule
  \end{tabular}
\end{table*}


\subsection{The Reference Forecast Changes Both Magnitude and Shape}
\label{sec:rq1-reference}

Holding the model, the data and the training procedure fixed, we
compute skill for the same forecasts against two references: smart
persistence and the convex reference. For gradient boosting on array
11 over daylight hours, skill against persistence is 0.252, 0.527 and
0.652 at one, three and six hours; against the convex reference, over
the same rows and the same forecasts, it is 0.200, 0.276 and 0.194
(Table~\ref{tab:T2}). At six hours, 70 percent of the apparent skill
is attributable to the reference. The pattern holds on all three
arrays: 0.244 to 0.190, 0.510 to 0.275 and 0.639 to 0.201 on array 12;
0.166 to 0.122, 0.341 to 0.183 and 0.386 to 0.132 on array 17.

The shape, not only the magnitude, changes. Measured against
persistence, skill rises monotonically with horizon - the pattern
commonly reported as evidence that a method's advantage grows at
longer lead times. Measured against the convex reference, skill is
non-monotonic and peaks at three hours, falling by six hours to
approximately its one-hour value. Fig.~\ref{fig:F1} shows this for all
three arrays: the solid curve (persistence) rises monotonically in
every panel; the dashed curve (convex reference) peaks at three hours
in every panel and falls back at six.

\begin{figure}[t]
  \centering
  \includegraphics[width=\columnwidth]{figures/F1_skill_vs_horizon.pdf}
  \caption{Forecast skill against two reference forecasts, XGBoost,
  lagged regime, validation year 2014, one panel per array. Skill
  against smart persistence (solid) rises monotonically with horizon
  in every panel; skill against the optimal convex combination of
  climatology and persistence (dashed) peaks at three hours and falls
  back at six. The apparent horizon trend is a property of the
  reference, not the model. Error bars show one standard deviation
  across five seeds and are smaller than the markers. Only XGBoost is
  shown; the other four architectures follow the same pattern.}
  \label{fig:F1}
\end{figure}

The mechanism is structural. Smart persistence forecasts $\hat{P}(t) =
k_p(t-h) \times P_{cs}(t)$, and at a six-hour horizon a midday target
is issued near dawn, when no valid clear-sky index yet exists: the
proportion of daylight forecasts relying on a forward-filled rather
than directly observed index rises from 3.3 percent at one hour to
22.6 percent at three and 51 percent at six, and the resulting bias
reaches $-2.15$ kW at midday on a 5 kW array. The fitted convex weight
$w$ reflects this directly, falling from 0.77 at one hour to 0.25 at
three and 0.04 at six on array 11 - at long horizons the optimal
reference is almost entirely climatology, which has no comparable
staleness problem. Array 17 is a partial exception: its fitted weight
remains higher at every horizon (0.83, 0.50, 0.31), because its
persistence forecast is stronger in relative terms, and its
non-monotonic peak in Fig.~\ref{fig:F1} is correspondingly shallower.

Under a paired Diebold-Mariano test with HAC (Bartlett-kernel)
variance, the Harvey-Leybourne-Newbold small-sample correction, and
Holm-adjusted p-values \cite{diebold1995comparing, harvey1997testing},
the convex reference beats smart persistence in all nine array-horizon
cells on both the validation and the test split ($p_{\text{holm}} <
0.05$ in all 18 comparisons): the reference this paper treats as
proper is not merely a construction, it is a significantly harder
target.

The shape inversion itself, not only the significance of the
reference gap, replicates on the held-out test year (2015), evaluated
once after every modelling and evaluation choice was frozen, on all
three arrays. Skill against persistence is 0.266, 0.574 and 0.701 on
array 11, 0.263, 0.556 and 0.684 on array 12, and 0.077, 0.281 and
0.381 on array 17, at one, three and six hours - monotonic throughout.
Against the convex reference it is 0.210, 0.316 and 0.232 on array 11,
0.202, 0.305 and 0.225 on array 12, and 0.043, 0.156 and 0.144 on
array 17: non-monotonic on every array, peaking at three hours and
falling back by six, on data no evaluation choice was tuned against.
Array 17's absolute skill shifts between years by more than the other
two arrays' does, and the shift traces to its reference rather than
its models: array 17's convex weight leans more heavily on
persistence at every horizon (0.31 to 0.83, against 0.04 to 0.77 for
arrays 11 and 12), so its reference is structurally more exposed to
persistence's year-to-year variation. This affects magnitude, not the
shape the paper claims; the full decomposition is in the anonymised
repository accompanying this paper.

\subsection{Night Hours Deflate Error, Not Skill}
\label{sec:rq1-night}

Restricting evaluation to daylight hours or retaining all 24 hours
changes reported error substantially and reported skill almost not at
all. If night-time errors are approximately zero, then
\begin{equation}
  \text{RMSE}_{\text{all}} = \text{RMSE}_{\text{daylight}} \times
  \sqrt{N_{\text{daylight}} / N_{\text{all}}}.
  \label{eq:closed-form}
\end{equation}
At six hours on array 11, $N_{\text{daylight}} = 3756$ and
$N_{\text{all}} = 8694$, giving a predicted ratio of 0.657 against an
observed 0.658; agreement is within 0.3 percent at every array and
horizon. Reported normalised error at this site is therefore deflated
by approximately 34 percent by a protocol choice unreported in 17 of
the 27 papers we surveyed (Table~\ref{tab:T1}), and the deflation
factor is a property of latitude and season, not of the forecasting
method. Skill, by contrast, is nearly invariant to the same choice -
for gradient boosting on array 11 it moves from 0.252 to 0.261, 0.527
to 0.528 and 0.652 to 0.652 at one, three and six hours - because the
reference forecast is deflated by the same factor as the model.

\subsection{Architecture Contributes Comparatively Little}
\label{sec:rq2-architecture}

Across five architectures - gradient boosting, a recurrent network, a
convolutional-recurrent network, and both recurrent variants with a
gradient-boosted residual correction stage - evaluated under the
identical protocol above, the largest difference we measure between
any two architectures at any array and horizon is 0.024 in skill
score, an order of magnitude smaller than either protocol effect
reported above.

\subsection{Held-Out Test Year}
\label{sec:test-split}

The reference-forecast shape inversion, shown above to hold on the
test split, is not the only direction that does. Night hours remain
deflationary and skill remains nearly unchanged by the same
closed-form argument, which does not depend on which year it is
applied to, and the largest architecture difference remains an order
of magnitude smaller than either protocol effect on the test split as
well as on validation.
